% from ega2-1-fr.tex
\oldpage[II]{46}

\begin{env}[2.8.12]
\label{II.2.8.12-fr}
Sous les hypothèses de (\hyperref[II.2.8.9-fr]{2.8.9}), pour
$f'\in S'_d$ ($d>0$) et $f=\varphi(f')$, le diagramme
\[
\label{II.2.8.12.1-fr}
  \xymatrix{
    M'_{(f')}\ar[r]^{\sim}\ar[d]
    & M'^{(d)}/(f'-1)M'^{(d)}\ar[d]\\
    M_{(f)}\ar[r]^{\sim}
    & M^{(d)}/(f-1)M^{(d)}
  }
\tag{2.8.12.1}
\]
(cf. (\hyperref[II.2.2.5-fr]{2.2.5})) est commutatif.
\end{env}

\begin{env}[2.8.13]
\label{II.2.8.13-fr}
Gardons les notations et hypothèses de (\hyperref[II.2.8.9-fr]{2.8.9}), et
soit $\mathcal{F}'$ un $\mathcal{O}_{X'}$-Module~; si on pose
$\mathcal{F}=\Phi^*(\mathcal{F}')$, on a, pour tout $n\in\mathbf{Z}$,
$\mathcal{F}(n)=\Phi^*(\mathcal{F}'(n))$ en vertu de
(\hyperref[II.2.8.11-fr]{2.8.11}) et
(\hyperref[0.4.3.3-fr]{0, 4.3.3}). Par suite
(\hyperref[0.3.7.1-fr]{0, 3.7.1}), on a un homomorphisme canonique
\[
  \Gamma(\rho):\Gamma(X',\mathcal{F}'(n))
  \longrightarrow\Gamma(X,\mathcal{F}(n))
\]
ce qui donne un di-homomorphisme canonique de modules gradués
\[
  \Gamma_\bullet(\mathcal{F}')\longrightarrow
  \Gamma_\bullet(\mathcal{F}).
\]

Supposons l'idéal $S_+$ engendré par $S_1$, et
$\mathcal{F}'=\widetilde{M'}$, donc $\mathcal{F}=\widetilde{M}$ avec
$M=M'\otimes_{A'}A$. Si $f'$ est homogène dans $S'_+$ et
$f=\varphi(f')$, on a vu que
$M_{(f)}=M'_{(f')}\otimes_{A'}A$, et le diagramme
\[
  \xymatrix{
    M'_0\ar[r]\ar[d]
    & M'_{(f')}\ar[d]
    & =\Gamma(D_+(f'),\widetilde{M'})\\
    M_0\ar[r]
    & M_{(f)}
    & =\Gamma(D_+(f),\widetilde{M})
  }
\]
est donc commutatif~; on conclut aussitôt de cette remarque et de la
définition de l'homomorphisme $\alpha$
(\hyperref[II.2.6.2-fr]{2.6.2}) que le diagramme
\[
\label{II.2.8.13.1-fr}
  \xymatrix{
    M'\ar[r]^{\alpha_{M'}}\ar[d]
    & \Gamma_\bullet(\widetilde{M'})\ar[d]\\
    M\ar[r]_{\alpha_M}
    & \Gamma_\bullet(\widetilde{M})
  }
\tag{2.8.13.1}
\]
est commutatif. De même, le diagramme
\[
\label{II.2.8.13.2-fr}
  \xymatrix{
    \widetilde{\Gamma_\bullet(\mathcal{F}')}\ar[r]^{\beta_{\mathcal{F}'}}
      \ar[d]
    & \mathcal{F}'\ar[d]\\
    \widetilde{\Gamma_\bullet(\mathcal{F})}\ar[r]_{\beta_{\mathcal{F}}}
    & \mathcal{F}
  }
\tag{2.8.13.2}
\]
est commutatif (la flèche verticale de droite étant le $\Phi$-morphisme
canonique
$\mathcal{F}'\to\Phi^*(\mathcal{F}')=\mathcal{F}$).
\end{env}

